\section{MLP 究竟是什么？参数、隐藏层和偏置}

先给出全称：\textbf{MLP = Multi-Layer Perceptron，多层感知机}。它并非一个神秘的新乘法，而是\textbf{线性变换、非线性激活、再线性变换}所组成的网络。

\subsection{从记忆检索写法到标准 MLP 写法}

前一节已经得到：
\[
\vect{y}=V^\top\relu(K\vect{x}).
\]
教材常用的两层 MLP 则写成：
\begin{equation}
\boxed{\vect{y}=W_2\sigma(W_1\vect{x}+\vect{b}_1)+\vect{b}_2.}
\end{equation}
这里 $W_1$ 可以对应前文的 $K$，$W_2$ 可以对应 $V^\top$。相比记忆检索版本，标准写法多出\textbf{偏置（bias）} $\vect{b}_1,\vect{b}_2$。

偏置先理解成乘完以后再加一个数。例如 $3x+2$ 比 $3x$ 多了一个常数 $2$；矩阵版本里这个常数变成一列数字。

\subsection{两个线性层怎样确定大小？}

设输入有 2 个数字，隐藏层有 3 个神经元，输出有 2 个数字：
\[
\underbrace{\vect{x}}_{2\times1}
\xrightarrow{\;W_1\;}
\underbrace{\vect{h}}_{3\times1}
\xrightarrow{\;\sigma\;}
\underbrace{\vect{a}}_{3\times1}
\xrightarrow{\;W_2\;}
\underbrace{\vect{y}}_{2\times1}.
\]
所以
\[
\underbrace{W_1}_{3\times2}\underbrace{\vect{x}}_{2\times1}
=\underbrace{\vect{h}}_{3\times1},
\qquad
\underbrace{W_2}_{2\times3}\underbrace{\vect{a}}_{3\times1}
=\underbrace{\vect{y}}_{2\times1}.
\]
这不是随便规定的：第一层希望输出 3 个数字，就需要 3 行；第二层从 3 个数字变成 2 个数字，就需要 2 行 3 列。

\subsection{如果中间没有非线性呢？}

忽略偏置且令 $\sigma$ 为恒等函数，得到
\[
\vect{y}=W_2(W_1\vect{x})=(W_2W_1)\vect{x}.
\]
两个线性层仍然等于\textbf{一个矩阵乘法}。有了 ReLU，不同输入会产生不同的零值模式，因此通常不能找到一个对所有输入都有效的固定矩阵替代这个网络。

\subsection{模型里真的有一行代表一个事实吗？}

未必。真实 MLP 的 $W_1,W_2$ 是通过大量数据与训练算法学得的。把 $W_1$ 的行理解成 Key、把 $W_2$ 的列理解成 Value，是一种很有用的\textbf{计算视角}，并不是说每一行权重必然编码了一条可以单独读出的事实。

\begin{conclusion}
MLP 的骨架就是：\textbf{第一次矩阵乘法产生隐藏层分数；激活函数逐元素改变分数；第二次矩阵乘法把它们组合为输出。}当你理解了矩阵的行与列，这个公式就不需要死背。
\end{conclusion}
